% !TeX root = ../prompt-engineering-guide-for-beginners.tex

%%%%%%%%%%%%%%%%%%%%%%%%%%%%%%%%%%
%% 第二章：动手试试
%%%%%%%%%%%%%%%%%%%%%%%%%%%%%%%%%%
\chapter{动手试试：五分钟开启你的第一次 AI 对话}
\label{ch:first-try}

上一章我们聊了 AI 是什么、不是什么。这一章，咱们直接动手。
读再多理论，不如自己试一次来得印象深刻。
五分钟之后，你就会完成人生中第一次（或者更有质量的一次）AI 对话。

\section{选一个工具，注册并打开}

首先，你需要选一个 AI 对话工具。如果你尚未注册过任何 AI 工具，推荐从以下两个对新手最友好的国产模型开始练习。它们的注册流程比海外产品简单，且交互界面十分友好：

\begin{itemize}
  \item \href{https://qwen.ai}{\textbf{通义千问 (Qwen)}}：由阿里巴巴推出，支持直接使用手机号注册，免费额度充足，且擅长处理中文的复杂语义。
  \item \href{https://chat.deepseek.com}{\textbf{DeepSeek}}：同样是优秀的国产工具，不仅注册简单无门槛，还在逻辑推理和代码生成上有亮眼表现，日常使用完全免费。
\end{itemize}

绝大多数平台的注册流程都非常类似，通常只需三步：

\begin{enumerate}
  \item 打开网站，点击「注册」或「Sign Up」。
  \item 填写邮箱或手机号，设置密码并完成验证。
  \item 登录后，你就能看到一个处于页面中央或底部的对话输入框。
\end{enumerate}

\begin{tipbox}
就如前面所说的，初学阶段不用纠结选哪个模型「最好」。工具只是载体，最重要的是掌握提示词的构建技巧。随意挑选一个顺眼的，现在就开始吧。
\end{tipbox}

\section{你的第一条提示词}

登录之后，你会看到一个输入框（通常在页面底部）。试试直接输入一句话：

\begin{promptbox}
帮我用三句话介绍杭州。
\end{promptbox}

按下回车（或发送按钮），等几秒钟，AI 就会给你一段回答。

你可能会得到类似这样的内容：

\begin{aibox}
杭州是浙江省的省会城市，以西湖闻名于世。这座城市有着两千多年的历史，是中国重要的电子商务中心。杭州也以龙井茶、丝绸和美食著称，是一座将古典韵味与现代活力完美融合的城市。
\end{aibox}

恭喜你，你已经完成了人生中的第一次 AI 对话。感觉是不是挺简单的（有手就会）？

不过你可能也会觉得：这个回答虽然没什么错误，但读起来像千篇一律的维基百科，没什么特别之处。
别着急，提示工程的魅力就在接下来的这一步。让我们试试换一种问法。

\section{同一个问题，换个问法试试}

现在，我们把同样的需求重新表达一遍，但这次加上具体的限制条件和细节要求：

\begin{promptbox}
你是一名旅游博主，请用轻松幽默的语气，用三句话向从没去过杭州的朋友介绍杭州，重点突出美食和自然风光。
\end{promptbox}

你可能会得到这样的回答：

\begin{aibox}
杭州是一座你去了就不想走的城市：西湖边散个步，感觉自己活成了古诗里的人。饿了就去河坊街，片儿川、叫花鸡、龙井虾仁轮番轰炸你的味蕾。如果你还没来过，那你的旅行清单上就少了最重要的一站。
\end{aibox}

对比一下两次回答，差别是不是很明显？

第一次的回答中规中矩，像百科全书；第二次活泼生动，像一个真正的旅游博主在跟你聊天。而我们做的唯一改变，就是在提示词里加了三样结构化约束条件：

\begin{enumerate}
  \item \textbf{角色设定（Role-playing）}：「你是一名旅游博主」
  \item \textbf{风格约束（Tone/Style）}：「轻松幽默的语气」
  \item \textbf{内容锚点（Focus/Content Anchors）}：「重点突出美食和自然风光」
\end{enumerate}

这就是提示词的「魔力」所在。当你没有给出任何约束时，模型会从它海量的训练数据中匹配最通用、最安全的回答方式。但当你给定了「角色」和「语气」等关键信息后，AI 就会把注意力集中在对应的风格和领域上。同一个 AI，能力没变，只是被你的提示词引导到了不同的方向，产出自然就不一样了。

\section{试试让 AI 帮你做一件真实的事}

理论讲完了，现在是你自己动手的时间。

请从你今天的待办清单里挑一件事，试着让 AI 帮你完成。比如：

\begin{itemize}
  \item 「帮我写一封给同事的邮件，通知下周三下午两点开项目复盘会议，语气正式但友好。」
  \item 「把下面这段英文翻译成中文，保持专业的语气：（粘贴你的英文内容）」
  \item 「帮我列一个周末两天的北京旅游计划，预算 2000 元以内，住宿和交通也帮我考虑进去。」
  \item 「我需要写一份季度工作总结，请帮我列一个大纲，包括工作成果、遇到的挑战和下季度计划。」
  \item 「我第一次请朋友来家里吃饭，帮我设计一桌 4 菜 1 汤的家常菜菜单，要有荤有素、好做又好看，买菜预算 100 元以内，并列出具体的食材采购清单。」
\end{itemize}

挑一个试试，真实地体验一下 AI 是如何帮你提高效率的。

如果它第一次给出的回答不够让你满意，完全没关系。你可以直接在对话框里追问和打磨，试着调整你的提示词：补充一些遗漏的细节、要求换一种说话方式，或者直接指令它「回答太冗长了，请精简并在 100 字以内总结要点」。

\begin{tipbox}
纸上得来终觉浅，绝知此事要躬行。亲自在对话框里敲下回车键，并观察 AI 随着你指令修改而发生的反馈变化，比阅读十页理论知识都要管用得多。
\end{tipbox}
